\begin{studiumworked}{M13. Respuesta lineal a un pequeño cambio de área · Intermedio}
\textbf{Enunciado.} El modelo normalizado $dP/d\theta=aP^n-P$ tiene $n=0.4$, $P(0)=1$. El área salta a $a=1.02$. Obtenga la respuesta lineal en $\theta=3$, el equilibrio lineal y el equilibrio no lineal.

\textbf{Datos.} Exponente: $0.4$; Escalón: $0.02$; Tiempo relativo: $3$.

\textbf{Incógnitas.} Magnitudes solicitadas y cierre físico.

\textbf{Hipótesis.} Modelo académico con hipótesis del capítulo. Datos originales; no representan hardware ni ensayo.

\textbf{Modelo físico.} Conservación y relaciones constitutivas declaradas.
\[P=1+\delta,\quad d\delta/d\theta=\epsilon-(1-n)\delta\]
\[\delta(\theta)=\epsilon[1-e^{-(1-n)\theta}]/(1-n)\]
\[P_{eq}=a^{1/(1-n)}\]
\textbf{Desarrollo.}
\begin{enumerate}
\item Linealizar alrededor de uno elimina productos de perturbaciones. El equilibrio lineal es $\delta_{eq}=0.02/0.6=0.03333333$.
\item A tiempo relativo tres, $\delta=0.02782337$ y $P\simeq1.027823$. El tiempo de relajación es $1/(1-n)=1.666667$ unidades de $\theta$.
\item El equilibrio no lineal tiene $P_{eq}-1=1.02^{1/0.6}-1=0.03355506$. La diferencia con $0.03333333$ es de orden cuadrático en el escalón.
\end{enumerate}
\textbf{Operaciones para contrastar el cálculo.}
\begin{enumerate}

\item Magnitud 1.

\[\fitmath{y_{1}=\frac{0.02}{0.6}\cdot \left(1-\exp\left(-0.6\cdot 3\right)\right)\approx 0.02782337\,}\]

\item Magnitud 2.

\[\fitmath{y_{2}={1.02}^{\frac{1}{0.6}}\approx 1.033555\,}\]

\end{enumerate}

\textbf{Resultados.} $P(3)\simeq1.027823$; $P_{eq,lineal}=1.033333$; $P_{eq,no\ lineal}=1.033555$.

\textbf{Interpretación.} El resultado corresponde a las hipótesis indicadas; su interpretación requiere conservar fronteras y unidades.

\textbf{Comprobación.} La respuesta arranca con delta cero y tiende al equilibrio lineal; el autovalor -0.6 es negativo.

\textbf{Límites.} Estabilidad local del modelo concentrado, no estabilidad acústica de un motor.
\end{studiumworked}
